\chapter{中心势与量子边界条件}
\label{chap:phys-hdqm-higher-dimensional-central-potentials}

三维中心势可用球谐函数分离变量，但把空间维数改成 $d$ 后，角向简并度、径向测度和离心项都会一起改变。若直接猜测三维公式的推广，很容易漏掉从 $r^{d-1}\dd r$ 到平直测度 $\dd r$ 的酉变换，并因此写错有效势。

在 \(d\) 维空间里，角变量由 \(SO(d)\) 的不可约表示组织，径向概率则带有体积因子 \(r^{d-1}\)。把超球谐本征值与这个测度一起处理，径向波函数可以酉地搬到 \(L^2(0,\infty;\dd r)\)；三维公式中熟悉的离心势随之变成一个依赖 \(d\) 和角量子数的逆平方项。这个项决定原点附近允许哪些局部解，也为下一章讨论算子定义域留下准确的边界数据。

\section{旋转对称、超球谐与径向自伴算子}
\label{sec:phys-hdqm-general-radial}

旋转群出现以前，先把三个名词分开。群是一批可以复合并可逆的操作；表示是把这些操作写成向量空间上的可逆线性算符；Lie 代数则只记录连续群在单位元附近的无穷小变化。它们不是同一个对象。

\begin{definitionplain}[旋转群 \(SO(d)\)]
\(d\) 维旋转群定义为
\[
SO(d)=\{R\in M_d(\mathbb R):R^{\mathsf T}R=I,\ \det R=1\}.
\]
条件 \(R^{\mathsf T}R=I\) 保证长度和夹角不变，\(\det R=1\) 排除镜面反射。矩阵乘法给出群乘法。
\end{definitionplain}

\begin{definitionplain}[群表示与生成元]
群 \(G\) 在向量空间 \(\mathcal H\) 上的表示，是满足
\(U(g_1g_2)=U(g_1)U(g_2)\) 和 \(U(e)=I\) 的映射
\(U:G\to GL(\mathcal H)\)。若 \(g(t)\) 是从单位元出发的一参数子群，则
\[
U(g(t))=\exp\!\left(-\frac{\ii t}{\hbar}J\right)
\]
中的自伴算符 \(J\) 称为这一变换的生成元。
\end{definitionplain}

\begin{definitionplain}[Casimir 算符]
若一个由生成元构成的算符 \(C\) 与表示中的每个生成元都对易，则称 \(C\) 为 Casimir 算符。它在每个不可约表示空间上是常数倍的恒等算符；这是 Schur 引理的直接结论。
\end{definitionplain}

先取整数 $d\ge2$；$d=1$ 的角向群是平凡的，后面在点相互作用与 Maxwell--Poisson 特例中单独处理。单粒子构型空间取 $\mathbb R^d$，Hilbert 空间为
\begin{equation}
\Hil=L^2(\mathbb R^d,\dd^d x),\qquad
H=-\frac{\hbar^2}{2\mu}\nabla_d^2+V(r),\qquad r=|\vb x|,
\label{eq:phys-hdqm-hilbert-hamiltonian}
\end{equation}
其中 $\mu$ 是约化质量，$V(r)$ 是实值中心势。旋转 $R\in SO(d)$ 在 $\Hil$ 上的酉表示定义为
\[
(U(R)\psi)(\vb x)=\psi(R^{-1}\vb x).
\]
对坐标平面 $(a,b)$ 的无穷小旋转，生成元为
\begin{equation}
L_{ab}=-\ii\hbar\left(x_a\partial_b-x_b\partial_a\right),
\qquad 1\le a<b\le d.
\label{eq:phys-hdqm-so-d-generators}
\end{equation}
并约定 $L_{ba}=-L_{ab}$、$L_{aa}=0$，从而可以在任意指标顺序下书写代数关系。它们满足 $\mathfrak{so}(d)$ 对易代数
\begin{equation}
[L_{ab},L_{cd}]
=\ii\hbar\left(
\delta_{ac}L_{bd}-\delta_{ad}L_{bc}
-\delta_{bc}L_{ad}+\delta_{bd}L_{ac}
\right).
\label{eq:phys-hdqm-so-d-algebra}
\end{equation}
中心势只依赖 $r$，所以 $[H,U(R)]=0$；等价地，$H$ 与所有 $L_{ab}$ 对易。角向谱因此应先由 $SO(d)$ 的二次 Casimir 组织，而不是从某个三维球谐公式类推。

定义
\begin{equation}
L^2\equiv\frac12\sum_{a,b=1}^dL_{ab}L_{ab}
=-\hbar^2\Delta_{S^{d-1}}.
\label{eq:phys-hdqm-casimir-sphere}
\end{equation}
单位球面 $S^{d-1}$ 上的超球谐函数 $Y_{l\alpha}(\Omega)$ 是这一 Casimir 的本征函数：
\begin{equation}
-\Delta_{S^{d-1}}Y_{l\alpha}
=l(l+d-2)Y_{l\alpha},
\qquad l=0,1,2,\ldots.
\label{eq:phys-hdqm-hyperspherical-spectrum}
\end{equation}
这里 $\alpha$ 标记固定 $l$ 的角向本征空间。对 $d\ge3$，其维数为
\begin{equation}
g_l^{(d)}=
\frac{(2l+d-2)(l+d-3)!}{l!(d-2)!},
\qquad d\ge3.
\label{eq:phys-hdqm-degeneracy}
\end{equation}
二维需要单独理解：$S^1$ 上的本征函数是 $\ee^{\ii m\phi}/\sqrt{2\pi}$，$m\in\mathbb Z$；若仍以 $l=|m|\ge0$ 分组，则 $g_0^{(2)}=1$，而 $g_l^{(2)}=2$ 对 $l\ge1$。因此下面以 $g_l^{(d)}$ 记固定 $l$ 的角向本征空间维数，但不把二维的 $\pm m$ 二重态误称为一个不可约 $SO(2)$ 表示。
取归一化
\[
\int_{S^{d-1}}Y_{l\alpha}^*(\Omega)Y_{l'\alpha'}(\Omega)\dd\Omega_{d-1}
=\delta_{ll'}\delta_{\alpha\alpha'},
\]
并记
\begin{equation}
S_{d-1}=\frac{2\pi^{d/2}}{\Gamma(d/2)},\qquad
\dd^d x=r^{d-1}\dd r\,\dd\Omega_{d-1}.
\label{eq:phys-hdqm-d-volume}
\end{equation}
于是 Hilbert 空间按角向 Laplace 本征空间分解为
\begin{equation}
\Hil\simeq
\bigoplus_{l=0}^{\infty}
\bigoplus_{\alpha=1}^{g_l^{(d)}}
L^2\!\left((0,\infty),r^{d-1}\dd r\right).
\label{eq:phys-hdqm-partial-wave-hilbert}
\end{equation}

\begin{proof}[交换关系]
先记 $D_{ab}=x_a\partial_b-x_b\partial_a$。基本恒等式
\[
[x_a\partial_b,x_c\partial_d]
=\delta_{bc}x_a\partial_d-\delta_{ad}x_c\partial_b
\]
逐项给出
\[
[D_{ab},D_{cd}]
=\delta_{bc}D_{ad}-\delta_{bd}D_{ac}
-\delta_{ac}D_{bd}+\delta_{ad}D_{bc}.
\]
代入 $L_{ab}=-\ii\hbar D_{ab}$，并用 $L_{ba}=-L_{ab}$ 重排指标，即得\eqnrefs{eq:phys-hdqm-so-d-algebra}。
\end{proof}

\begin{proof}[Casimir 与球面 Laplace 算子]
令 $D=r\partial_r=\sum_ax_a\partial_a$。直接展开生成元的平方，利用
\[
\sum_{a,b}x_a^2\partial_b^2=r^2\nabla_d^2,
\qquad
\sum_{a,b}x_ax_b\partial_a\partial_b=D^2-D,
\]
以及
\[
\sum_{a,b}x_a\delta_{ab}\partial_b=D,
\]
得到
\[
\frac12\sum_{a,b}(x_a\partial_b-x_b\partial_a)^2
=r^2\nabla_d^2-D^2-(d-2)D.
\]
另一方面，球坐标中的 Laplace 算子满足
\[
r^2\nabla_d^2=D^2+(d-2)D+\Delta_{S^{d-1}}.
\]
两式相减即得
\[
\frac12\sum_{a,b}(x_a\partial_b-x_b\partial_a)^2
=\Delta_{S^{d-1}}.
\]
再乘上 $(-\ii\hbar)^2=-\hbar^2$，便得到\eqnrefs{eq:phys-hdqm-casimir-sphere}。
\end{proof}

\begin{proof}[超球谐本征值与简并度]
取次数为 $l$ 的齐次多项式 $P_l(\vb x)$，写成
\[
P_l(\vb x)=r^lY_l(\Omega).
\]
把球坐标 Laplace 算子作用到它上面，得到
\[
\begin{aligned}
\nabla_d^2[r^lY_l]
&=\frac1{r^{d-1}}\partial_r
\left(r^{d-1}\partial_r r^l\right)Y_l
+\frac{r^l}{r^2}\Delta_{S^{d-1}}Y_l\\
&=r^{l-2}
\left[l(l+d-2)+\Delta_{S^{d-1}}\right]Y_l.
\end{aligned}
\]
若 $P_l$ 是调和齐次多项式，即 $\nabla_d^2P_l=0$，便有
\[
-\Delta_{S^{d-1}}Y_l=l(l+d-2)Y_l,
\]
这就推出\eqnrefs{eq:phys-hdqm-hyperspherical-spectrum}。

对 $d\ge3$，次数 $l$ 的全部齐次多项式空间维数为
\[
\dim\mathcal P_l=\binom{l+d-1}{l}.
\]
Fischer 分解
\[
\mathcal P_l=\mathcal H_l\oplus r^2\mathcal P_{l-2}
\]
把它分成调和子空间 $\mathcal H_l$ 与带一个 $r^2$ 因子的低两次部分；当 $l<2$ 时约定 $\mathcal P_{l-2}=\{0\}$，等价地下面第二个二项式系数取零。因此
\[
\begin{aligned}
g_l^{(d)}
&=\dim\mathcal H_l
=\binom{l+d-1}{l}-\binom{l+d-3}{l-2}\\
&=\frac{(2l+d-2)(l+d-3)!}{l!(d-2)!},
\end{aligned}
\]
即得\eqnrefs{eq:phys-hdqm-degeneracy}。二维 $l=0$ 的阶乘边界情形应单独处理，而 $l\ge1$ 时同一维数差给出 $2$，对应 $m=\pm l$。因此 $l(l+d-2)$ 与 $g_l^{(d)}$ 都来自 $SO(d)$ 作用下调和多项式的表示结构，而不是对三维球谐公式的形式替换。
\end{proof}

例：三维 $d=3$ 有 $g_l=2l+1$（$m=-l,\ldots,l$），即标准球谐函数 $Y_l^m$；二维 $d=2$ 有 $g_0=1$、$g_l=2$（$l\ge1$，对应 $\cos l\theta$ 和 $\sin l\theta$），即 Fourier 模；四维 $d=4$ 有 $g_l=(l+1)^2$。

对固定 $(l,\alpha)$ 写
\[
\psi(\vb x)=R_l(r)Y_{l\alpha}(\Omega).
\]
由
\[
\nabla_d^2=\frac{1}{r^{d-1}}\frac{\dd}{\dd r}
\left(r^{d-1}\frac{\dd}{\dd r}\right)
+\frac1{r^2}\Delta_{S^{d-1}},
\]
得到径向 Schr\"odinger 方程
\begin{equation}
R_l''+\frac{d-1}{r}R_l'
+\left[
\frac{2\mu}{\hbar^2}(E-V(r))
-\frac{l(l+d-2)}{r^2}
\right]R_l=0.
\label{eq:phys-hdqm-radial-r}
\end{equation}
定义
\begin{equation}
\nu_l=l+\frac{d-2}{2},\qquad
l=0,1,2,\ldots,
\label{eq:phys-hdqm-nu}
\end{equation}
并作酉变换
\[
u_l(r)=r^{(d-1)/2}R_l(r).
\]
这里左边的 $u_l$ 是约化径向波函数，不应与阶数 $\nu_l$ 混淆。归一化满足
\[
\int_0^\infty |R_l(r)|^2r^{d-1}\dd r
=\int_0^\infty |u_l(r)|^2\dd r,
\]
故径向问题被搬到 $L^2(0,\infty;\dd r)$。方程化为
\begin{equation}
-u_l''+
\left[
\frac{\nu_l^2-\frac14}{r^2}
+\frac{2\mu}{\hbar^2}V(r)
\right]u_l
=\frac{2\mu E}{\hbar^2}u_l.
\label{eq:phys-hdqm-radial-u}
\end{equation}
因此应先定义最小对称算子
\begin{equation}
 h_l^{\min}
=-\frac{\dd^2}{\dd r^2}
+\frac{\nu_l^2-\frac14}{r^2}
+\frac{2\mu}{\hbar^2}V(r),
\qquad
\Dom(h_l^{\min})=C_c^\infty(0,\infty).
\label{eq:phys-hdqm-radial-minimal-operator}
\end{equation}
物理哈密顿量是 $h_l^{\min}$ 的一个自伴实现。特殊函数只能求解相应微分表达式；量子动力学是否唯一，还要由定义域和端点边界条件决定。

\begin{proof}[解]
 令 \(q=(d-1)/2\)，作代换 \(R_l=r^{-q}u_l\) 消去径向拉普拉斯算子的一阶导数项，使方程化为标准 Sturm--Liouville 形式，再定出逆平方系数与原点端点分类。

 逐次求导，
 \begin{equation*}
  R_l'=r^{-q}\left(u_l'-\frac qru_l\right),
  \qquad
  R_l''=r^{-q}\left[
   u_l''-\frac{2q}{r}u_l'
   +\frac{q(q+1)}{r^2}u_l\right].
 \end{equation*}
 与一阶径向项合并，利用 $d-1=2q$，$u_l'$ 项相消，得
 \begin{equation*}
  R_l''+\frac{d-1}{r}R_l'
  =r^{-q}\left[
   u_l''-\frac{(d-1)(d-3)}{4r^2}u_l
  \right].
 \end{equation*}
 代回\eqnrefs{eq:phys-hdqm-radial-r}，逆平方系数变为
 \begin{equation*}
  \begin{aligned}
  l(l+d-2)+\frac{(d-1)(d-3)}4
  &=l^2+(d-2)l+\frac{(d-1)(d-3)}4\\
  &=\left(l+\frac{d-2}{2}\right)^2-\frac14\\
  &=\nu_l^2-\frac14,
  \end{aligned}
 \end{equation*}
 其中 \(\nu_l:=l+(d-2)/2\)。\(\nu_l\) 即原点附近两个独立解的幂指数差，决定波函数的原点奇性与自伴性。两个独立解在原点附近的行为为
 \begin{equation*}
   u_l^{(1)}(r)\sim r^{\nu_l+1/2},\qquad
   u_l^{(2)}(r)\sim r^{-\nu_l+1/2},
 \end{equation*}
 当 \(\nu_l\ge1\) 时 \(u_l^{(2)}\) 不平方可积，原点为极限点型；当 \(0\le\nu_l<1\) 时两支皆平方可积，原点为极限圆型。于是得到\eqnrefs{eq:phys-hdqm-radial-u}。同时
 \begin{equation*}
  \|R_l\|_{L^2(r^{d-1}\dd r)}^2
  =\int_0^\infty r^{d-1}r^{-(d-1)}|u_l|^2\dd r
  =\|u_l\|_{L^2(\dd r)}^2,
 \end{equation*}
 故 $R_l\mapsto u_l=r^{(d-1)/2}R_l$ 确为径向 Hilbert 空间之间的酉映射。

 常用维数中 \(d=3\) 给出 \(u_l=rR_l\)、\(\nu_l=l+1/2\) 与离心项 \(l(l+1)/r^2\)；\(d=2\) 给出 \(\nu_l=l\)，其 \(s\) 波恰处系数 \(-1/4\) 的临界逆平方势。三维与二维的 \(s\) 波均属极限圆型，通常的正则性条件只是相应自伴实现中的一个选择。
 \end{proof}

对最大算子定义域中的 $u,v$，分部积分给出 Green 边界形式
\begin{equation}
\langle u,h_lv\rangle-\langle h_lu,v\rangle
=\left[(u')^*v-u^*v'\right]_{0}^{\infty}.
\label{eq:phys-hdqm-boundary-form}
\end{equation}
若无穷远为极限点型，唯一可能的扩张自由度来自 $r=0$。设原点附近
\[
V(r)=\frac{v_{-2}}{r^2}+o(r^{-2}),
\]
则总逆平方系数为
\begin{equation}
g_l^{\rm eff}
=\nu_l^2-\frac14+\frac{2\mu v_{-2}}{\hbar^2},
\qquad
\chi_l=\sqrt{g_l^{\rm eff}+\frac14}.
\label{eq:phys-hdqm-effective-inverse-square}
\end{equation}
当 $g_l^{\rm eff}>-1/4$ 时，两个 Frobenius 分支为
\[
u(r)\sim A r^{1/2-\chi_l}+B r^{1/2+\chi_l}.
\]
两支在原点都平方可积当且仅当 $0\le\chi_l<1$。于是
\begin{equation}
\begin{cases}
\chi_l\ge1\;(g_l^{\rm eff}\ge3/4): & r=0\text{ 为 limit-point，原点不需要额外边界条件},\\
0\le\chi_l<1\;(-1/4\le g_l^{\rm eff}<3/4): & r=0\text{ 为 limit-circle，需要一参数自伴边界条件}.
\end{cases}
\label{eq:phys-hdqm-limit-point-circle}
\end{equation}
对 $0<\chi_l<1$，把渐近系数记为 $(A[u],B[u])$。代入\eqnrefs{eq:phys-hdqm-boundary-form} 可得
\[
\lim_{r\to0}[(u')^*v-u^*v']
=2\chi_l\left(B[u]^*A[v]-A[u]^*B[v]\right).
\]
因此任一实参数 $\Lambda_l\in\mathbb R\cup\{\infty\}$ 所规定的
\begin{equation}
B[u]=\Lambda_l A[u]
\label{eq:phys-hdqm-extension-boundary}
\end{equation}
都使边界形式在定义域内消失；这里 $\Lambda_l=\infty$ 表示 $A[u]=0$，而有限 $\Lambda_l$ 的量纲为 $\mathrm{length}^{-2\chi_l}$。这就是自伴扩张参数在波函数原点渐近式中的具体位置。

这里的极限圆型自由度属于定义在半直线 $(0,\infty)$ 上、原点已被去掉的最小径向算子。若原始问题确实定义在完整的 $\mathbb R^d$ 上，且原点没有额外的点缺陷或比 $r^{-2}$ 更奇异的短程相互作用，则完整哈密顿量的定义域还要满足原点的通常正则性；这会从上述扩张族中选出一个规范实现，而不是额外引入新的物理参数。例如 $d=3,l=0$ 且 $v_{-2}=0$ 时 $\chi_l=\nu_l=1/2$，两支为 $u(r)\sim A+B r$。由于 $R(r)=u(r)/r$，完整三维空间上的通常正则（等价于该半有界问题的 Friedrichs）实现要求 $A=0$，即在\eqnrefs{eq:phys-hdqm-extension-boundary} 的记号下选择 $\Lambda_l=\infty$。二维 $s$ 波的临界情形则有 $R(r)\sim A+B\ln(r/L_0)$，完整平面上的无缺陷拉普拉斯算子选择 $B=0$。其余允许的扩张正是给原点增加零程边界数据；它们将在后面的点相互作用一节中作为新的物理模型讨论。

临界点 $g_l^{\rm eff}=-1/4$ 时两根合并，
\[
u(r)\sim\sqrt r\left[A+B\ln\frac r{L_0}\right],
\]
自伴边界条件固定 $A/B$ 的一个实值，等价地引入一个长度尺度。若 $g_l^{\rm eff}<-1/4$，写 $\chi_l=\ii s_l$，则
\[
u(r)\sim\sqrt r\left[
C\ee^{+\ii s_l\ln(r/L_0)}
+D\ee^{-\ii s_l\ln(r/L_0)}
\right].
\]
自伴性要求入射与出射的对数波具有相同模，即 $D=\ee^{\ii\theta_l}C$。改变参考尺度 $L_0$ 只会平移 $\theta_l$，故这一相位等价于引入一个短程尺度；但纯粹选择自伴扩张并不能恢复半有界性：超临界 $-1/r^2$ 哈密顿量仍有向 $-\infty$ 延伸的几何束缚态塔。稳定的物理模型必须由有限尺度的短程完成截断这一紫外坍缩。

\begin{proof}[解]
考虑原点主导方程
\begin{equation*}
 -u''+\frac g{r^2}u=0.
\end{equation*}
设 $u=r^\alpha$，则
\begin{equation*}
 u'=\alpha r^{\alpha-1},
 \qquad
 u''=\alpha(\alpha-1)r^{\alpha-2}.
\end{equation*}
代入方程并除以 $r^{\alpha-2}$，得到指标方程
\begin{equation*}
 -\alpha(\alpha-1)+g=0.
\end{equation*}
改写为
\begin{equation*}
 \alpha^2-\alpha-g=0,
 \qquad
 \alpha=\frac{1\pm\sqrt{1+4g}}2
 =\frac12\pm\sqrt{g+\frac14}.
\end{equation*}
当 $g>-1/4$ 时记 $\chi=\sqrt{g+1/4}$，两支为
\begin{equation*}
 u_\pm(r)=r^{1/2\pm\chi}.
\end{equation*}
约化径向波函数的局部平方可积要求 $\int_0^\epsilon|u|^2\dd r<\infty$，于是
\begin{equation*}
 \int_0^\epsilon r^{2\alpha_\pm}\dd r
 =\int_0^\epsilon r^{1\pm2\chi}\dd r.
\end{equation*}
对 $+$ 分支，被积函数幂次 $1+2\chi>0$ 恒成立，积分总有限。对 $-$ 分支，要求
\begin{equation*}
 1-2\chi>-1
 \quad\Longleftrightarrow\quad
 \chi<1
 \quad\Longleftrightarrow\quad
 g+\frac14<1
 \quad\Longleftrightarrow\quad
 g<\frac34.
\end{equation*}
因此当 $g\ge3/4$ 时 $-$ 分支不再局部平方可积，定义域只剩一个 $L^2$ 解，为极限点型；当 $-1/4\le g<3/4$ 时两支皆可积，为极限圆型，需要额外边界条件才能选定自伴扩张。临界 $g=-1/4$ 时两根合并 $\alpha=1/2$，第二解由降次法得到 $u_2=\sqrt r\ln r$，仍属极限圆型。当 $g<-1/4$ 时记 $s=\sqrt{-g-1/4}$，两根为
\begin{equation*}
 \alpha=\frac12\pm\ii s,
 \qquad
 u_\pm(r)=r^{1/2}\ee^{\pm\ii s\ln r},
\end{equation*}
两支模长都为 $r^{1/2}$，皆局部平方可积；但能量二次型
\begin{equation*}
 \int_0^\epsilon\!\left(|u'|^2+\frac g{r^2}|u|^2\right)\dd r
\end{equation*}
在 $g<-1/4$ 时失去下界，进入 fall-to-the-中心区域，自伴扩张虽仍存在却无法恢复半有界性。
\end{proof}
